\chapter{Projectvoorstel}
\label{app:projectvoorstel}

\section{Projectcontext}
Vesalius.ai ontwikkelt software die ziekenhuizen en artsen ondersteunt bij
administratieve processen en patiëntcommunicatie. Binnen het platform worden
automatisch e-mails verstuurd op vooraf bepaalde momenten, bijvoorbeeld bij
het bevestigen, herinneren of annuleren van een afspraak.

De inhoud en opmaak van deze e-mails zijn grotendeels vooraf bepaald.
Organisaties kunnen bepaalde visuele elementen, zoals het logo en de
merkkleuren, configureren, maar hebben niet voor elk e-mailtype de mogelijkheid
om de inhoud en opmaak zelfstandig aan te passen.

Dit project onderzoekt hoe een visuele e-maileditor deze beperking kan helpen
oplossen.

\section{Probleemstelling}
Gebruikers hebben onvoldoende vrijheid om de inhoud en opmaak van bepaalde
automatisch gegenereerde e-mails aan te passen. Wanneer wijzigingen aan de
e-mailinhoud nodig zijn, is hiervoor mogelijk technische tussenkomst vereist.

Daarnaast moet de gegenereerde HTML correct worden verwerkt en weergegeven
in verschillende e-mailclients. Een visueel correcte weergave in de editor
garandeert immers niet automatisch dat de e-mail er in elke e-mailclient
hetzelfde uitziet.

De centrale onderzoeksvraag luidt:

\begin{quote}
Hoe kan een veilige en gebruiksvriendelijke WYSIWYG-e-maileditor worden
ontwikkeld waarmee gebruikers e-mailtemplates visueel kunnen aanpassen en
waarvan de gegenereerde e-mails betrouwbaar kunnen worden getest in gangbare
e-mailclients?
\end{quote}

\section{Doelstelling}
Het doel is een zelfstandige proof of concept te ontwikkelen waarmee
gebruikers de inhoud en basisopmaak van vooraf bepaalde e-mailtemplates
visueel kunnen aanpassen, opslaan, vooraf bekijken en testen.

De PoC moet aantonen dat de editor kan worden gecombineerd met een
renderingproces dat HTML verwerkt, toegestane placeholders invult en de
e-mail voorbereidt voor verzending en weergave in verschillende
e-mailclients.

De PoC vormt een technische basis voor een mogelijke latere integratie in
het Vesalius.ai-platform. Een dergelijke integratie valt buiten de scope
van dit project.

\section{Functionele vereisten}
De PoC moet de volgende functionaliteiten ondersteunen:

\begin{enumerate}
    \item Een bestaande e-mailtemplate selecteren en laden.
    \item De inhoud van een template visueel bewerken met een WYSIWYG-editor.
    \item Ondersteunde tekstopmaak toepassen.
    \item Wijzigingen opslaan en opnieuw laden.
    \item Een voorbeeld van de gegenereerde e-mail bekijken.
    \item Toegestane placeholders met fictieve gegevens invullen.
    \item De gegenereerde HTML veilig verwerken.
    \item Een testmail versturen naar een lokale testomgeving.
    \item De werking en uitvoer controleren met geautomatiseerde tests.
\end{enumerate}

\section{Afbakening}

\subsection{Binnen de scope}

\begin{itemize}
    \item Een zelfstandige frontend met een visuele e-maileditor.
    \item Een mock-API voor het laden en opslaan van templates.
    \item Een renderingpipeline voor HTML, placeholders en opmaak.
    \item Een previewfunctie en een lokale testmailomgeving.
    \item Drie piloottemplates: afspraakbevestiging, afspraakherinnering en
    afspraakannulering.
    \item Geautomatiseerde tests en compatibiliteitscontroles.
    \item Documentatie van de technische keuzes, beperkingen en resultaten.
\end{itemize}

\subsection{Buiten de scope}

\begin{itemize}
    \item Integratie in de productieomgeving van Vesalius.ai.
    \item Het effectief versturen of plannen van productie-e-mails.
    \item Het aanpassen van Keycloak-e-mails.
    \item Een volledige editor voor vrij opgebouwde e-maillay-outs.
    \item Het verwerken van echte patiëntgegevens.
    \item Het ontwikkelen van een afzonderlijke AI-functionaliteit.
\end{itemize}

\section{Technologische aanpak}
De beoogde architectuur bestaat uit een Angular-frontend en een afzonderlijke
mock-API op basis van Node.js, TypeScript en Express.

Voor de editor wordt Quill 2 als kandidaat onderzocht. De uiteindelijke keuze
wordt onderbouwd met een technische proef. Voor de renderingpipeline worden
onder meer HTML-sanitatie, een gecontroleerde placeholderlijst en het inline
plaatsen van CSS onderzocht.

De testomgeving kan gebruikmaken van Docker en Mailpit. Voor unit- en
integratietests wordt Vitest overwogen en voor browsergebaseerde tests
Playwright. De uiteindelijke technologiekeuzes en versies worden vastgelegd
in het hoofdstuk over technologie.

\section{Verwachte resultaten}
Aan het einde van het project worden de volgende resultaten nagestreefd:

\begin{itemize}
    \item Een werkende, zelfstandige PoC van de visuele e-maileditor.
    \item Een werkende laad-, opslag- en renderingflow voor templates.
    \item Een preview- en testmailfunctionaliteit.
    \item Een testset met fictieve gegevens en geautomatiseerde controles.
    \item Een evaluatie van de weergave in beschikbare e-mailclients.
    \item Technische documentatie en een voorstel voor mogelijke latere
    integratie.
\end{itemize}

Dit zijn beoogde resultaten. De uiteindelijke realisatie en eventuele
beperkingen worden beschreven in het hoofdstuk over de resultaten.

\section{Privacy en veiligheid}
Tijdens de ontwikkeling worden uitsluitend fictieve gegevens gebruikt.
HTML-inhoud wordt volgens een expliciet toegestane lijst verwerkt en
placeholders worden gecontroleerd voordat ze worden ingevuld.

Productiegegevens, productiegeheimen en echte patiëntinformatie maken geen
deel uit van de PoC. De oplossing wordt niet beschouwd als een productieklare
e-mailinfrastructuur.

\section{Opvolging en evaluatie}
De voortgang wordt opgevolgd aan de hand van de sprintplanning en de
bijbehorende taken. Per sprint worden de uitgevoerde werkzaamheden,
opgeleverde resultaten, eventuele problemen en volgende stappen geregistreerd.

De uiteindelijke evaluatie gebeurt aan de hand van de functionele vereisten,
de testresultaten, de renderingkwaliteit en de vastgestelde beperkingen.